\begin{problem}[理论考试][转速测量与控制装置]{25}
  一转速测量和控制装置的原理如图所示。在 O 点有电量为 $Q$ 的正电荷，内壁光滑的轻质绝缘细管可绕通过 O 点的竖直轴在水平面内转动，在管内距离 O 为 $L$ 处有一光电触发控制开关 A，在 O 端固定有一自由长度为 $L/4$ 的轻质绝缘弹簧，弹簧另一端与一质量为 $m$、带有正电荷 $q$ 的小球相连接。

  开始时，系统处于静态平衡。细管在外力矩作用下，作定轴转动，小球可在细管内运动。当细管转速 $\omega$ 逐渐变大时，小球到达细管的A处刚好相对于细管径向平衡，并触发控制开关，外力矩瞬时变为零，从而限制转速过大；同时O点的电荷变为等量负电荷 $-Q$。通过测量此后小球相对于细管径向平衡点的位置B，可测定转速。若测得OB的距离为 $L/2$，求
  \begin{subproblem}
    弹簧系数 $k_{0}$ 及小球在 B 处时细管的转速；
  \end{subproblem}
  \begin{subproblem}
    试问小球在平衡点 B 附近是否存在相对于细管的径向微振动？如果存在，求出该微振动的周期。
  \end{subproblem}
\end{problem}